%!TeX root=../princesstop.tex
\chapter{»I Tried Not To Be«}

\begin{figure}[t!]
\centering
\includegraphics[width=\linewidth]{filigree}
\end{figure}

\lettrine[lines=5]{I}{t} was pretty, comfortable Mrs Carmichael who explained everything. She was sent for at once, and came across the square to take Sara into her warm arms and make clear to her all that had happened. The excitement of the totally unexpected discovery had been temporarily almost overpowering to Mr Carrisford in his weak condition.

»Upon my word,« he said faintly to Mr Carmichael, when it was suggested that the little girl should go into another room. »I feel as if I do not want to lose sight of her.«

»I will take care of her,« Janet said, »and mamma will come in a few minutes.« And it was Janet who led her away.

»We're so glad you are found,« she said. »You don't know how glad we are that you are found.«

Donald stood with his hands in his pockets, and gazed at Sara with reflecting and self-reproachful eyes.

»If I'd just asked what your name was when I gave you my sixpence,« he said, »you would have told me it was Sara Crewe, and then you would have been found in a minute.« Then Mrs Carmichael came in. She looked very much moved, and suddenly took Sara in her arms and kissed her.

»You look bewildered, poor child,« she said. »And it is not to be wondered at.«

Sara could only think of one thing.

»Was he,« she said, with a glance toward the closed door of the library—»was \textsc{he} the wicked friend? Oh, do tell me!«

Mrs Carmichael was crying as she kissed her again. She felt as if she ought to be kissed very often because she had not been kissed for so long.

»He was not wicked, my dear,« she answered. »He did not really lose your papa's money. He only thought he had lost it; and because he loved him so much his grief made him so ill that for a time he was not in his right mind. He almost died of brain fever, and long before he began to recover your poor papa was dead.«

»And he did not know where to find me,« murmured Sara. »And I was so near.« Somehow, she could not forget that she had been so near.

»He believed you were in school in France,« Mrs Carmichael explained. »And he was continually misled by false clues. He has looked for you everywhere. When he saw you pass by, looking so sad and neglected, he did not dream that you were his friend's poor child; but because you were a little girl, too, he was sorry for you, and wanted to make you happier. And he told Ram Dass to climb into your attic window and try to make you comfortable.«

Sara gave a start of joy; her whole look changed.

»Did Ram Dass bring the things?« she cried out. »Did he tell Ram Dass to do it? Did he make the dream that came true?«

»Yes, my dear—yes! He is kind and good, and he was sorry for you, for little lost Sara Crewe's sake.«

The library door opened and Mr Carmichael appeared, calling Sara to him with a gesture.

»Mr Carrisford is better already,« he said. »He wants you to come to him.«

Sara did not wait. When the Indian gentleman looked at her as she entered, he saw that her face was all alight.

She went and stood before his chair, with her hands clasped together against her breast.

»You sent the things to me,« she said, in a joyful emotional little voice, »the beautiful, beautiful things? \textsc{you} sent them!«

»Yes, poor, dear child, I did,« he answered her. He was weak and broken with long illness and trouble, but he looked at her with the look she remembered in her father's eyes—that look of loving her and wanting to take her in his arms. It made her kneel down by him, just as she used to kneel by her father when they were the dearest friends and lovers in the world.

»Then it is you who are my friend,« she said; »it is you who are my friend!« And she dropped her face on his thin hand and kissed it again and again.

»The man will be himself again in three weeks,« Mr Carmichael said aside to his wife. »Look at his face already.«

In fact, he did look changed. Here was the »Little Missus,« and he had new things to think of and plan for already. In the first place, there was Miss Minchin. She must be interviewed and told of the change which had taken place in the fortunes of her pupil.

Sara was not to return to the seminary at all. The Indian gentleman was very determined upon that point. She must remain where she was, and Mr Carmichael should go and see Miss Minchin himself.

»I am glad I need not go back,« said Sara. »She will be very angry. She does not like me; though perhaps it is my fault, because I do not like her.«

But, oddly enough, Miss Minchin made it unnecessary for Mr Carmichael to go to her, by actually coming in search of her pupil herself. She had wanted Sara for something, and on inquiry had heard an astonishing thing. One of the housemaids had seen her steal out of the area with something hidden under her cloak, and had also seen her go up the steps of the next door and enter the house.

»What does she mean!« cried Miss Minchin to Miss Amelia.

»I don't know, I'm sure, sister,« answered Miss Amelia. »Unless she has made friends with him because he has lived in India.«

»It would be just like her to thrust herself upon him and try to gain his sympathies in some such impertinent fashion,« said Miss Minchin. »She must have been in the house for two hours. I will not allow such presumption. I shall go and inquire into the matter, and apologize for her intrusion.«

Sara was sitting on a footstool close to Mr Carrisford's knee, and listening to some of the many things he felt it necessary to try to explain to her, when Ram Dass announced the visitor's arrival.

Sara rose involuntarily, and became rather pale; but Mr Carrisford saw that she stood quietly, and showed none of the ordinary signs of child terror.

Miss Minchin entered the room with a sternly dignified manner. She was correctly and well dressed, and rigidly polite.

»I am sorry to disturb Mr Carrisford,« she said; »but I have explanations to make. I am Miss Minchin, the proprietress of the Young Ladies' Seminary next door.«

The Indian gentleman looked at her for a moment in silent scrutiny. He was a man who had naturally a rather hot temper, and he did not wish it to get too much the better of him.

»So you are Miss Minchin?« he said.

»I am, sir.«

»In that case,« the Indian gentleman replied, »you have arrived at the right time. My solicitor, Mr Carmichael, was just on the point of going to see you.«

Mr Carmichael bowed slightly, and Miss Minchin looked from him to Mr Carrisford in amazement.

»Your solicitor!« she said. »I do not understand. I have come here as a matter of duty. I have just discovered that you have been intruded upon through the forwardness of one of my pupils—a charity pupil. I came to explain that she intruded without my knowledge.« She turned upon Sara. »Go home at once,« she commanded indignantly. »You shall be severely punished. Go home at once.«

The Indian gentleman drew Sara to his side and patted her hand.

»She is not going.«

Miss Minchin felt rather as if she must be losing her senses.

»Not going!« she repeated.

»No,« said Mr Carrisford. »She is not going home—if you give your house that name. Her home for the future will be with me.«

Miss Minchin fell back in amazed indignation.

»With \textsc{you}! With \textsc{you,} sir! What does this mean?«

»Kindly explain the matter, Carmichael,« said the Indian gentleman; »and get it over as quickly as possible.« And he made Sara sit down again, and held her hands in his—which was another trick of her papa's.

Then Mr Carmichael explained—in the quiet, level-toned, steady manner of a man who knew his subject, and all its legal significance, which was a thing Miss Minchin understood as a business woman, and did not enjoy.

»Mr Carrisford, madam,« he said, »was an intimate friend of the late Captain Crewe. He was his partner in certain large investments. The fortune which Captain Crewe supposed he had lost has been recovered, and is now in Mr Carrisford's hands.«

»The fortune!« cried Miss Minchin; and she really lost colour as she uttered the exclamation. »Sara's fortune!«

»It \textsc{will} be Sara's fortune,« replied Mr Carmichael, rather coldly. »It is Sara's fortune now, in fact. Certain events have increased it enormously. The diamond mines have retrieved themselves.«

»The diamond mines!« Miss Minchin gasped out. If this was true, nothing so horrible, she felt, had ever happened to her since she was born.

»The diamond mines,« Mr Carmichael repeated, and he could not help adding, with a rather sly, unlawyer-like smile, »There are not many princesses, Miss Minchin, who are richer than your little charity pupil, Sara Crewe, will be. Mr Carrisford has been searching for her for nearly two years; he has found her at last, and he will keep her.«

After which he asked Miss Minchin to sit down while he explained matters to her fully, and went into such detail as was necessary to make it quite clear to her that Sara's future was an assured one, and that what had seemed to be lost was to be restored to her tenfold; also, that she had in Mr Carrisford a guardian as well as a friend.

Miss Minchin was not a clever woman, and in her excitement she was silly enough to make one desperate effort to regain what she could not help seeing she had lost through her worldly folly.

»He found her under my care,« she protested. »I have done everything for her. But for me she should have starved in the streets.«

Here the Indian gentleman lost his temper.

»As to starving in the streets,« he said, »she might have starved more comfortably there than in your attic.«

»Captain Crewe left her in my charge,« Miss Minchin argued. »She must return to it until she is of age. She can be a parlour boarder again. She must finish her education. The law will interfere in my behalf.«

»Come, come, Miss Minchin,« Mr Carmichael interposed, »the law will do nothing of the sort. If Sara herself wishes to return to you, I dare say Mr Carrisford might not refuse to allow it. But that rests with Sara.«

»Then,« said Miss Minchin, »I appeal to Sara. I have not spoiled you, perhaps,« she said awkwardly to the little girl; »but you know that your papa was pleased with your progress. And—ahem—I have always been fond of you.«

Sara's green-gray eyes fixed themselves on her with the quiet, clear look Miss Minchin particularly disliked.

»Have \textsc{you}, Miss Minchin?« she said. »I did not know that.«

Miss Minchin reddened and drew herself up.

»You ought to have known it,« said she; »but children, unfortunately, never know what is best for them. Amelia and I always said you were the cleverest child in the school. Will you not do your duty to your poor papa and come home with me?«

Sara took a step toward her and stood still. She was thinking of the day when she had been told that she belonged to nobody, and was in danger of being turned into the street; she was thinking of the cold, hungry hours she had spent alone with Emily and Melchisedec in the attic. She looked Miss Minchin steadily in the face.

»You know why I will not go home with you, Miss Minchin,« she said; »you know quite well.«

A hot flush showed itself on Miss Minchin's hard, angry face.

»You will never see your companions again,« she began. »I will see that Ermengarde and Lottie are kept away\longdash«

Mr Carmichael stopped her with polite firmness.

»Excuse me,« he said; »she will see anyone she wishes to see. The parents of Miss Crewe's fellow-pupils are not likely to refuse her invitations to visit her at her guardian's house. Mr Carrisford will attend to that.«

It must be confessed that even Miss Minchin flinched. This was worse than the eccentric bachelor uncle who might have a peppery temper and be easily offended at the treatment of his niece. A woman of sordid mind could easily believe that most people would not refuse to allow their children to remain friends with a little heiress of diamond mines. And if Mr Carrisford chose to tell certain of her patrons how unhappy Sara Crewe had been made, many unpleasant things might happen.

»You have not undertaken an easy charge,« she said to the Indian gentleman, as she turned to leave the room; »you will discover that very soon. The child is neither truthful nor grateful. I suppose«—to Sara—»that you feel now that you are a princess again.«

Sara looked down and flushed a little, because she thought her pet fancy might not be easy for strangers—even nice ones—to understand at first.

»I—\textsc{tried} not to be anything else,« she answered in a low voice—»even when I was coldest and hungriest—I tried not to be.«

»Now it will not be necessary to try,« said Miss Minchin, acidly, as Ram Dass salaamed her out of the room.


She returned home and, going to her sitting room, sent at once for Miss Amelia. She sat closeted with her all the rest of the afternoon, and it must be admitted that poor Miss Amelia passed through more than one bad quarter of an hour. She shed a good many tears, and mopped her eyes a good deal. One of her unfortunate remarks almost caused her sister to snap her head entirely off, but it resulted in an unusual manner.

»I'm not as clever as you, sister,« she said, »and I am always afraid to say things to you for fear of making you angry. Perhaps if I were not so timid it would be better for the school and for both of us. I must say I've often thought it would have been better if you had been less severe on Sara Crewe, and had seen that she was decently dressed and more comfortable. I \textsc{know} she was worked too hard for a child of her age, and I know she was only half fed\longdash«

»How dare you say such a thing!« exclaimed Miss Minchin.

»I don't know how I dare,« Miss Amelia answered, with a kind of reckless courage; »but now I've begun I may as well finish, whatever happens to me. The child was a clever child and a good child—and she would have paid you for any kindness you had shown her. But you didn't show her any. The fact was, she was too clever for you, and you always disliked her for that reason. She used to see through us both\longdash«

»Amelia!« gasped her infuriated elder, looking as if she would box her ears and knock her cap off, as she had often done to Becky.

But Miss Amelia's disappointment had made her hysterical enough not to care what occurred next.

»She did! She did!« she cried. »She saw through us both. She saw that you were a hard-hearted, worldly woman, and that I was a weak fool, and that we were both of us vulgar and mean enough to grovel on our knees for her money, and behave ill to her because it was taken from her—though she behaved herself like a little princess even when she was a beggar. She did—she did—like a little princess!« And her hysterics got the better of the poor woman, and she began to laugh and cry both at once, and rock herself backward and forward.

»And now you've lost her,« she cried wildly; »and some other school will get her and her money; and if she were like any other child she'd tell how she's been treated, and all our pupils would be taken away and we should be ruined. And it serves us right; but it serves you right more than it does me, for you are a hard woman, Maria Minchin, you're a hard, selfish, worldly woman!«

And she was in danger of making so much noise with her hysterical chokes and gurgles that her sister was obliged to go to her and apply salts and sal volatile to quiet her, instead of pouring forth her indignation at her audacity.

And from that time forward, it may be mentioned, the elder Miss Minchin actually began to stand a little in awe of a sister who, while she looked so foolish, was evidently not quite so foolish as she looked, and might, consequently, break out and speak truths people did not want to hear.

That evening, when the pupils were gathered together before the fire in the schoolroom, as was their custom before going to bed, Ermengarde came in with a letter in her hand and a queer expression on her round face. It was queer because, while it was an expression of delighted excitement, it was combined with such amazement as seemed to belong to a kind of shock just received.

»What \textsc{is} the matter?« cried two or three voices at once.

»Is it anything to do with the row that has been going on?« said Lavinia, eagerly. »There has been such a row in Miss Minchin's room, Miss Amelia has had something like hysterics and has had to go to bed.«

Ermengarde answered them slowly as if she were half stunned.

»I have just had this letter from Sara,« she said, holding it out to let them see what a long letter it was.

»From Sara!« Every voice joined in that exclamation.

»Where is she?« almost shrieked Jessie.

»Next door,« said Ermengarde, »with the Indian gentleman.«

»Where? Where? Has she been sent away? Does Miss Minchin know? Was the row about that? Why did she write? Tell us! Tell us!«

There was a perfect babel, and Lottie began to cry plaintively.

Ermengarde answered them slowly as if she were half plunged out into what, at the moment, seemed the most important and self-explaining thing.

»There \textsc{were} diamond mines,« she said stoutly; »there \textsc{were}!« Open mouths and open eyes confronted her.

»They were real,« she hurried on. »It was all a mistake about them. Something happened for a time, and Mr Carrisford thought they were ruined\longdash«

»Who is Mr Carrisford?« shouted Jessie.

»The Indian gentleman. And Captain Crewe thought so, too—and he died; and Mr Carrisford had brain fever and ran away, and \textsc{he} almost died. And he did not know where Sara was. And it turned out that there were millions and millions of diamonds in the mines; and half of them belong to Sara; and they belonged to her when she was living in the attic with no one but Melchisedec for a friend, and the cook ordering her about. And Mr Carrisford found her this afternoon, and he has got her in his home—and she will never come back—and she will be more a princess than she ever was—a hundred and fifty thousand times more. And I am going to see her tomorrow afternoon. There!«

Even Miss Minchin herself could scarcely have controlled the uproar after this; and though she heard the noise, she did not try. She was not in the mood to face anything more than she was facing in her room, while Miss Amelia was weeping in bed. She knew that the news had penetrated the walls in some mysterious manner, and that every servant and every child would go to bed talking about it.

So until almost midnight the entire seminary, realizing somehow that all rules were laid aside, crowded round Ermengarde in the schoolroom and heard read and re-read the letter containing a story which was quite as wonderful as any Sara herself had ever invented, and which had the amazing charm of having happened to Sara herself and the mystic Indian gentleman in the very next house.

Becky, who had heard it also, managed to creep up stairs earlier than usual. She wanted to get away from people and go and look at the little magic room once more. She did not know what would happen to it. It was not likely that it would be left to Miss Minchin. It would be taken away, and the attic would be bare and empty again. Glad as she was for Sara's sake, she went up the last flight of stairs with a lump in her throat and tears blurring her sight. There would be no fire tonight, and no rosy lamp; no supper, and no princess sitting in the glow reading or telling stories—no princess!

She choked down a sob as she pushed the attic door open, and then she broke into a low cry.

The lamp was flushing the room, the fire was blazing, the supper was waiting; and Ram Dass was standing smiling into her startled face.

»Missee sahib remembered,« he said. »She told the sahib all. She wished you to know the good fortune which has befallen her. Behold a letter on the tray. She has written. She did not wish that you should go to sleep unhappy. The sahib commands you to come to him tomorrow. You are to be the attendant of missee sahib. Tonight I take these things back over the roof.«

And having said this with a beaming face, he made a little salaam and slipped through the skylight with an agile silentness of movement which showed Becky how easily he had done it before.